%%% variant: default
\section{Forward \& Backward}

\begin{itemize}
\item Aus Text können wir Trainingspaare bilden: Eingabetoken → richtiges nächstes Token.
\begin{itemize}
\item Mit dem richtigen nächsten Token als Ziel können wir den Fehler berechnen.
\item Backpropagation optimiert sowohl die Netzwerkgewichte als auch die Embeddings.
\end{itemize}
\end{itemize}
%%% variant: instructor
\section{Forward \& Backward}

\begin{itemize}
\item Aus Text können wir Trainingspaare bilden: Eingabetoken → richtiges nächstes Token.
\begin{itemize}
\item Mit dem richtigen nächsten Token als Ziel können wir den Fehler berechnen.
\item Backpropagation optimiert sowohl die Netzwerkgewichte als auch die Embeddings.
\end{itemize}
\end{itemize}
%%% variant: handout
\section{Forward \& Backward}

\begin{itemize}
\item Aus Text können wir Trainingspaare bilden: Eingabetoken → richtiges nächstes Token.
\begin{itemize}
\item Mit dem richtigen nächsten Token als Ziel können wir den Fehler berechnen.
\item Backpropagation optimiert sowohl die Netzwerkgewichte als auch die Embeddings.
\end{itemize}
\end{itemize}
%%% variant: slides
\begin{frame}[t]
\frametitle{Forward \& Backward}

\end{frame}

\begin{frame}[t]
\begin{itemize}
\item Aus Text können wir Trainingspaare bilden: Eingabetoken → richtiges nächstes Token.
\begin{itemize}
\item Mit dem richtigen nächsten Token als Ziel können wir den Fehler berechnen.
\item Backpropagation optimiert sowohl die Netzwerkgewichte als auch die Embeddings.
\end{itemize}
\end{itemize}
\end{frame}

%%% variant: slides-dozent-h1-frame
\begin{frame}[t]
\frametitle{Forward \& Backward}

\begin{itemize}
\item Aus Text können wir Trainingspaare bilden: Eingabetoken → richtiges nächstes Token.
\begin{itemize}
\item Mit dem richtigen nächsten Token als Ziel können wir den Fehler berechnen.
\item Backpropagation optimiert sowohl die Netzwerkgewichte als auch die Embeddings.
\end{itemize}
\end{itemize}
\end{frame}

%%% variant: slides-h1-hide-level-1
\begin{frame}[t]

\begin{itemize}
\item Aus Text können wir Trainingspaare bilden: Eingabetoken → richtiges nächstes Token.
\begin{itemize}
\item Mit dem richtigen nächsten Token als Ziel können wir den Fehler berechnen.
\item Backpropagation optimiert sowohl die Netzwerkgewichte als auch die Embeddings.
\end{itemize}
\end{itemize}
\end{frame}

